\section{Introduction}

The advent of Large Language Models (LLMs) unlocks unprecedented capabilities in natural language understanding and generation, fundamentally transforming how machines process and generate human language. Recent breakthroughs in multimodal large language models~\cite{Bai2025Qwen3VL, comanici2025gemini_gemini25, gpt5, claudeopus45, guo2025seed1_seed15} further extend these capabilities to visual perception and action planning, providing the foundational abilities necessary for GUI automation. These advances enable models to understand complex visual interfaces and reason about appropriate interactions, opening new possibilities for autonomous GUI agents that can assist users in their daily digital tasks.

Building on these foundational capabilities, the research community makes significant progress in GUI agent development. Various approaches explore methods for acquiring GUI interaction data~\cite{gao2024mobileviews, rawles2023androidinthewild, liu2025scalecua, wang2025opencua}, establishing critical data infrastructure for training GUI agents. Several powerful GUI agent models~\cite{qin2025ui_uitars1, wang2025ui_uitars2, ye2025mobile_mobileagentv3, zeng2025uitron} demonstrate impressive capabilities in navigating and manipulating graphical interfaces. However, a critical challenge remains: \textit{how to efficiently obtain high-quality trajectory and knowledge data to improve agent performance in target domains?} Traditional annotation methods suffer from subjectivity and prohibitive costs, limiting the scalability of GUI agent development.

To address this challenge, we introduce a self-evolving training pipeline centered on the Calibrated Step Reward System (CSRS). While prior work explores self-improvement through data generation via rationalization, self-critique, quality filtering, and offline reinforcement learning~\cite{zelikman2022star, madaan2023self_selfrefine, jung2023impossibledistill, gulcehre2023reinforced_rest}, these approaches primarily target single-turn tasks where model-generated annotations can be reliably verified. 
For multi-turn GUI agents where task execution spans multiple steps, direct model-generated annotations are prone to factual errors and hallucinations in the absence of verifiable ground truth. 
Unlike existing methods that rely solely on model introspection or post-hoc filtering, CSRS anchors LLM-generated dense step-level reasoning to trajectory-level evaluation signals, either automated verification scripts or human annotations, thereby ensuring annotation quality while maintaining scalability. 
Through trajectory-level calibration and LLM-powered knowledge extraction, CSRS converts model-generated trajectories into high-quality training data, achieving $>$90\% annotation accuracy with 10-100$\times$ cost reduction compared to traditional step-level annotation. Our progressive three-stage training paradigm orchestrates parallel data flows for novel exploration and strategic knowledge filtering, continuously improving model capabilities across multiple training rounds. 

On top of this pipeline, we train \textbf{Step-GUI}, a family of GUI specialist models (4B, 8B) built upon the Qwen3-VL backbone, and Step-GUI-8B achieves state-of-the-art performance: 80.2\% on AndroidWorld, 48.5\% on OSWorld, 65\% on ScreenShot-Pro, and 89.91\%/52.50\% on our AndroidDaily benchmark (static/end-to-end). Remarkably, our compact 4B model delivers competitive performance while remaining deployable on consumer-grade hardware, enabling local execution without cloud dependencies.

As GUI agents gain enhanced capabilities in visual understanding and autonomous task execution, two fundamental challenges emerge: \textit{standardized communication across heterogeneous devices} and \textit{user privacy protection when processing sensitive data}. The fragmentation of development efforts and high integration costs historically suppress innovation in tool ecosystems~\cite{qu2025toollearning}. To address the standardization challenge, the Model Context Protocol (MCP)~\cite{mcp} is proposed as a universal standard that defines how agents and tools should communicate, enabling an interoperable ``plug-and-play'' ecosystem. Recent work explores how agents can effectively retrieve and select appropriate MCP services from large pools~\cite{gan2025rag_ragmcp} and how to convert existing software projects into MCP format~\cite{ouyang2025code2mcp}. Building on these foundations, we propose \textbf{GUI-MCP} (Graphical User Interface - Model Context Protocol), the first MCP implementation specifically designed for GUI automation that addresses both standardization and privacy protection. GUI-MCP provides a standardized, cross-platform protocol that abstracts device capabilities into a small set of atomic and composite tools. Its hierarchical dual-layer architecture combines a \textit{Low-level MCP} that offers fine-grained operations (e.g., clicks, swipes, text input) with a \textit{High-level MCP} that delegates entire tasks to locally deployed GUI specialist models such as Step-GUI-4B. This design allows the main LLM to focus on high-level planning while it offloads routine GUI operations to local models. Critically, GUI-MCP supports a high-privacy execution mode in which raw screenshots and sensitive states stay on the device and only semantic summaries flow to external LLMs, which effectively protects user privacy while it leverages cloud-based reasoning capabilities.

As the community makes progress on these technical challenges, attention is shifting toward a more pressing concern: \textit{how to evaluate whether GUI agents can handle real-world everyday usage beyond synthetic setups?} Existing work develops diverse evaluation benchmarks, ranging from static benchmarks~\cite{cheng2024seeclick, wu2024atlas, li2025screenspotpro, li2025autogui, burns2022dataset_motif, rastogi2021uibert_refexp, zhang2024llamatouch, liu2024visualwebbench_vwb} that assess element grounding and task planning through prediction-annotation matching, to interactive benchmarks~\cite{rawles2024androidworld, xie2024osworld} that enable task execution in realistic operating system environments. However, these efforts either concentrate on single-step prediction accuracy or lack coverage of high-frequency Chinese mobile applications, leaving a critical gap: \textit{can agents reliably handle the high-frequency, everyday tasks that constitute real-world mobile usage?} To address this gap, we introduce \textbf{AndroidDaily}, a benchmark explicitly grounded in empirical analysis of authentic mobile usage patterns. Rather than pursuing maximal application coverage, AndroidDaily focuses on ubiquitous daily scenarios (Transportation, Shopping, Social Media, Entertainment, Local Services) where agent deployment has immediate practical impact. The benchmark employs a two-tier evaluation strategy: a \textit{Static Benchmark} with 3146 actions for efficient single-step action prediction, and an \textit{End-to-End Benchmark} with 235 tasks that span multiple dimensions (scenario, task type, complexity, and ambiguity) and evaluate autonomous task completion in fully functional environments. Step-GUI-8B demonstrates strong performance on AndroidDaily, which highlights both its current capabilities and remaining challenges in realistic deployment scenarios.



In summary, our contributions establish a comprehensive framework for building practical GUI agents:
\begin{itemize}
\item A self-evolving training pipeline with CSRS that achieves 10-100$\times$ cost reduction while it continuously improves model capabilities through closed-loop data refinement.
\item Step-GUI models that achieve state-of-the-art performance across multiple benchmarks (8B: 80.2\% on AndroidWorld, 48.5\% on OSWorld, 62.6\% on ScreenShot-Pro, 89.91\%/ 52.50\% on AndroidDaily static/end-to-end), with the 4B variant that enables consumer-grade local deployment.
\item GUI-MCP, the first standardized protocol for LLM-device interaction that balances execution efficiency with privacy protection through hierarchical architecture.
\item AndroidDaily, an ecologically valid benchmark grounded in real-world usage patterns with dual evaluation modes (3146 static actions, 235 end-to-end tasks) and multi-dimensional taxonomies for targeted capability analysis.
\end{itemize}

Together, these contributions address the complete pipeline from model training to standardized deployment interfaces and authentic evaluation, which paves the way for GUI agents that can assist users in their daily mobile interactions.